\section{Introduction}
\label{sec:intro}

The scaling of deep foundation models has transformed artificial intelligence, enabling unprecedented capabilities in computer vision and natural language processing \cite{radford2021learning,touvron2023llama,devlin2018bert,vaswani2017attention}. However, training these massive parameterized architectures from scratch or performing full joint multi-task fine-tuning is exceptionally resource-intensive, requiring thousands of GPU hours and massive, centralized multi-domain datasets. Consequently, practitioners typically fine-tune pre-trained models on specialized individual tasks, resulting in multiple independent models that are highly capable in their respective domains but suffer from catastrophic forgetting on other tasks \cite{saim2026merge,wortsman2022robust}.

To consolidate these distinct capabilities without the prohibitive overhead of retraining, model merging has emerged as a compelling, zero-training-cost alternative \cite{ilharco2022editing,yadav2023resolving,wortsman2022model}. By operating directly in the weight space, methods like Task Arithmetic \cite{ilharco2022editing}, Ties-Merging \cite{yadav2023resolving}, DARE \cite{yu2023language}, SyMerge \cite{kim2024symerge}, and OrthoMerge \cite{spherelab2024orthomerge} successfully blend parameter updates (task vectors) to build unified, multi-task systems.

\begin{figure}[t]
\centering
\includegraphics[width=\columnwidth]{task_vectors.png}
\caption{\textbf{Overview of QP-Merge.} Rather than performing standard model merging and direct uniform quantization (which stretches quantization scales and ruins low-bit precision), QP-Merge decomposes task vectors into a dense quantized path (INT4/INT8) and a sparse high-precision outlier path (FP16). Layer-wise scale calibration is optimized over a tiny unlabeled calibration set to ensure near-lossless reconstruction.}
\label{fig:overview}
\end{figure}

Despite the elegance of these methods, academic literature almost exclusively evaluates model merging under high-precision floating-point formats (FP16 or FP32). In real-world enterprise applications and edge deployments---such as mobile phones, automotive systems, and IoT devices---serving large floating-point models is highly impractical. Instead, models must be compressed via Post-Training Quantization (PTQ) to low-bit integers, such as INT4 or INT8, to meet tight VRAM bandwidth limits, reduce serving latency, and minimize computational energy costs \cite{han2015deep,dettmers2022gpt3,xiao2023smoothquant,nagel2020up}. 

When standard post-training quantization is applied to merged multi-task models, we observe a severe and catastrophic degradation in task performance. For instance, directly quantizing a merged Vision Transformer (ViT-B-32) model to INT4 causes the average classification accuracy to drop by \textbf{3.42\%}, with SVHN classification dropping by a devastating \textbf{6.40\%} (from 90.72\% down to 84.32\%). This failure is driven by two key physical phenomena:
\begin{enumerate}
    \item \textbf{Heavy-tailed task outliers:} Task vector subtraction ($\Delta W_t = W_t - W_{\text{base}}$) inherently highlights highly sparse, high-magnitude parameter differences representing localized specialized features. In standard uniform quantization, these extreme weight outliers stretch the symmetric quantization scale factor $S_c$, compressing the dense majority of normal-range weights into only a handful of quantization bins, causing massive quantization distortion \cite{squeezellm2023}.
    \item \textbf{Activation scale mismatches:} Because the underlying task models are fine-tuned on entirely different distributions (e.g., MNIST vs. SVHN), their internal intermediate activation ranges and scale distributions are fundamentally mismatched. Linearly blending task vectors under a fixed global coefficient ignores these task-specific activation scale discrepancies, leading to severe representation misalignment after quantization.
\end{enumerate}

To bridge the gap between academic model merging and practical low-cost edge deployment, we propose \textbf{QP-Merge} (Quantization-Preserving Merging), a training-free framework that co-designs parameter merging and post-training quantization. Guided by a pragmatic research philosophy, QP-Merge introduces two lightweight, highly effective techniques:
\begin{itemize}
    \item \textbf{Outlier-Residual Decoupling (ORD):} We identify and isolate the top $\le 1\%$ extreme task-specific weight outliers from each task vector and route them to a high-precision, highly sparse tensor ($W_{\text{outlier}}$) represented in FP16 format. This leaves a tightly bounded, outlier-free dense remainder ($W_{\text{dense}}$) that can be aggressively quantized to 4-bit or 8-bit integers without stretching quantization scales.
    \item \textbf{Quantization-Error Aware Scale Calibration (QE-Calib):} Over a tiny calibration set of only $M=128$ completely unlabeled target domain samples, we perform a rapid, 100-step optimization of layer-wise diagonal weight scaling parameters $D_l$ and merging coefficients $\lambda$ to minimize activation reconstruction mean-squared error (MSE). This calibration aligns mismatched task scale dynamics directly in the quantized state.
\end{itemize}

Empirical evaluations on dual-task vision classification using a pre-trained `ViT-B-32` demonstrate that QP-Merge completely resolves the low-bit quantization bottleneck. In 4-bit (INT4) mode, QP-Merge achieves an average accuracy of \textbf{94.70\%} across 3 random seeds, recovering over \textbf{88\%} of the performance drop caused by naive quantization relative to the optimized baseline and performing \textbf{within 0.42\%} of the optimized unquantized FP32 merged bound. In 8-bit (INT8) mode, QP-Merge achieves an average accuracy of \textbf{95.14\%}, representing a \textbf{+0.02\%} gain over the optimized unquantized FP32 baseline (and a \textbf{+0.21\%} gain over the uniform unquantized FP32 reference), demonstrating that unsupervised joint scaling and merging coefficient calibration directly optimizes parameter fusion in quantized regimes. Crucially, QP-Merge is hardware-friendly, compatible with standard sparse matrix-multiply (SpMM) runtimes, and requires zero labeled downstream data, making it highly robust and instantly deployable in production settings.
